\documentclass[11pt,a4paper]{article}
\usepackage[T1]{fontenc}
\usepackage{lmodern}
\usepackage[margin=26mm]{geometry}
\usepackage{amsmath,amssymb,amsthm,mathtools}
\usepackage{microtype,booktabs,array}
\usepackage{xcolor}
\usepackage{xurl}
\usepackage[colorlinks=true,linkcolor=blue!45!black,citecolor=blue!45!black,urlcolor=blue!45!black]{hyperref}
\hypersetup{pdftitle={Escaping Sequences under Theta-Source Orthogonality Constraints},pdfauthor={Jongmin Choi},pdfsubject={A research note on constrained escape for a minimum theta-weighted jump form},pdfkeywords={theta source, closed quadratic form, orthogonality constraints, divided differences}}
\newcommand{\R}{\mathbb R}
\newcommand{\C}{\mathbb C}
\newcommand{\HH}{\mathcal H}
\newcommand{\DD}{\mathcal D}
\newcommand{\EE}{\mathcal E}
\newcommand{\CC}{\mathcal C}
\newcommand{\NN}{\mathcal N}
\newcommand{\Id}{I}
\newcommand{\ip}[2]{\langle #1,#2\rangle}
\newcommand{\norm}[1]{\lVert #1\rVert}
\newcommand{\abs}[1]{\lvert #1\rvert}
\newcommand{\dd}{\,\mathrm d}
\newcommand{\Var}{\operatorname{Var}}
\newcommand{\Dom}{\operatorname{D}}
\newcommand{\supp}{\operatorname{supp}}
\newcommand{\sech}{\operatorname{sech}}
\newcommand{\Etheta}{\mathbb E_{\theta}}
\newtheorem{theorem}{Theorem}[section]
\newtheorem{proposition}[theorem]{Proposition}
\newtheorem{lemma}[theorem]{Lemma}
\newtheorem{corollary}[theorem]{Corollary}
\theoremstyle{definition}
\newtheorem{definition}[theorem]{Definition}
\theoremstyle{remark}
\newtheorem{remark}[theorem]{Remark}
\numberwithin{equation}{section}
\setlength{\emergencystretch}{2em}
\allowdisplaybreaks[1]
\newcommand{\AuthorEmail}{24ping@naver.com}
\newcommand{\AuthorORCID}{0009-0008-7448-514X}
\title{Escaping Sequences under\\Theta-Source Orthogonality Constraints}
\author{Jongmin Choi\\[2pt]
\small Independent Researcher, Seoul, Korea\\
\small Email: \href{mailto:\AuthorEmail}{\texttt{\AuthorEmail}}\\
\small ORCID: \href{https://orcid.org/\AuthorORCID}{\AuthorORCID}}
\date{Research note --- 17 October 2026}

\begin{document}
\maketitle
\begin{abstract}
We consider the minimum closed even jump form associated with the infinite
theta source, an archimedean translation kernel, and all prime-power shifts.
Let $\NN$ be the closed space generated by the constant and the centered
even derivative ratios of the source. We construct real smooth unit vectors
$h_m\in\DD\cap\NN^\perp$ for which
$\abs{\EE(h_m)-1/2}\le C m^3e^{-\sqrt m}$ and the mass in every fixed
bounded interval tends to zero. Every defining orthogonality condition is
exact. The construction uses confluent divided differences at clusters of
zeros of the completed zeta function and the unconditional zero-counting
asymptotic; neither real nor simple zeros are assumed. We derive the source
identity from the classical compact-test explicit formula and justify all
minimum-domain admissions. For a positive two-shift envelope we also prove
a ground-state identity and show that its transformed energy tends to zero
on the same sequence. These results obstruct a positive uniform margin and
fixed-center coercivity under the known source constraints. They do not
determine the critical sign of the form, identify the full nullspace, or
give an operator Weyl sequence. No conclusion about the Riemann hypothesis
is obtained.
\end{abstract}

\section{The form and the constrained escape statement}\label{sec:setting}

For $x\ge0$ define
\begin{equation}\label{eq:kappa}
 \kappa(x)=e^{x/2}\sum_{n\ge1}(4z_n(x)^2-6z_n(x))e^{-z_n(x)},
 \qquad z_n(x)=\pi n^2e^{2x},
\end{equation}
and extend $\kappa$ evenly. Set
\begin{equation}\label{eq:weight}
 c(x)=\cosh(x/2),\qquad
 \dd\nu(x)=2\kappa(x)c(x)\dd x,\qquad
 \HH=L^2_{\mathrm{even}}(\R,\nu;\C).
\end{equation}
The normalization $\nu(\R)=1$ is proved below. All inner products are
conjugate-linear in the first variable. We use
$\tau_su(x)=u(x+s)$ and
$\Var_\nu(h)=\norm h_{\HH}^2-\abs{\int h\dd\nu}^2$. Write
\begin{equation}\label{eq:measure}
 \begin{split}
 \rho(s)&=\frac{e^{-s/2}}{1-e^{-2s}},\quad s>0,\qquad
 c_n=\frac{\Lambda(n)}{\sqrt n},\quad s_n=\log n,\\
 \dd\mu(s)&=\frac12\rho(\abs{s})\dd s
       +\frac12\sum_{n\ge2}c_n(\delta_{s_n}+\delta_{-s_n}).
 \end{split}
\end{equation}
Here $\Lambda$ includes every prime power. The positive form on
$\CC=C_c^\infty(\R;\C)\cap\HH$ is
\begin{align}\label{eq:energy}
 \EE(h)
 &=\iint\kappa(x)\kappa(z)\abs{h(x)-h(z)}^2
                 \dd\mu(x-z)\dd x\notag\\
 &=\int_0^\infty\rho(s)\int_\R\kappa(x)\kappa(x+s)
                   \abs{h(x+s)-h(x)}^2\dd x\dd s\notag\\
 &\quad+\sum_{n\ge2}c_n\int_\R\kappa(x)\kappa(x+s_n)
                   \abs{h(x+s_n)-h(x)}^2\dd x.
\end{align}
Thus the first line uses a symmetric half measure. In a double integral
with $\dd\mu(x-z)\dd x$, the variable $z$ is integrated against the translated
shift measure; atoms are not replaced by a Lebesgue density.

The form is closable, and $(\EE,\DD)$ denotes its \emph{minimum} closure
from $\CC$ in the norm
\begin{equation}\label{eq:form-norm}
 \norm h_{\DD}^2=\EE(h)+\norm h_{\HH}^2.
\end{equation}
No equality with a maximal jump-energy domain is assumed. Let $A\ge0$ be
the representing self-adjoint operator and set
\begin{equation}\label{eq:L}
 P_1h=\ip1h_{\HH}\,1,\qquad
 L=A-\frac12I+\frac12P_1,\qquad \Dom(L)=\Dom(A).
\end{equation}
The known source-generated space considered here is
\begin{equation}\label{eq:N}
 f_j=\frac{\kappa^{(2j)}}{\kappa}\quad(j\ge0),\qquad
 \NN=\overline{\operatorname{span}}_{\HH}\{f_j:j\ge0\}.
\end{equation}
We prove $f_j\in\HH$ and $\ip1{f_j}_{\HH}=4^{-j}$ below, so this is also
the closed span of $1$ and $f_j-4^{-j}$ for $j\ge1$.

\begin{theorem}[Escape under every known source constraint]\label{thm:escape}
There exist constants $C<\infty$, $m_0$, and real smooth even functions
$h_m\in\DD\cap\NN^\perp$ for all integers $m\ge m_0$ such that
\begin{equation}\label{eq:main-energy}
 \norm{h_m}_{\HH}=1,\qquad
 \abs{\EE(h_m)-\tfrac12}\le C m^3e^{-\sqrt m}.
\end{equation}
For every fixed $T<\infty$,
\begin{equation}\label{eq:main-local}
 \norm{\mathbf1_{[-T,T]}h_m}_{\HH}\longrightarrow0,
 \qquad h_m\rightharpoonup0\text{ in }\HH.
\end{equation}
Each moment $\ip{f_j}{h_m}_{\HH}$ is exactly zero for every $j\ge0$.
For each fixed admissible envelope $R$ of Definition~\ref{def:envelope},
there are constants $C_R<\infty$ and $a_R>0$ such that, after increasing
the starting index if necessary,
\begin{equation}\label{eq:main-ground}
 0\le D_R(h_m)\le C_R m^3e^{-a_R\sqrt m}\longrightarrow0.
\end{equation}
\end{theorem}

\begin{corollary}\label{cor:obstructions}
On $\DD\cap\NN^\perp$ there is no inequality
$\EE(h)\ge(1/2+\varepsilon)\norm h_{\HH}^2$ with $\varepsilon>0$.
For fixed finite $T$ there is no $a>0$ such that
$\norm{\mathbf1_{[-T,T]}h}_{\HH}\ge a\norm h_{\HH}$, even among unit
vectors with $\EE(h)\le1/2+\varepsilon$ for any fixed $\varepsilon>0$.
For an admissible fixed $R$, no finite nonnegative $a,b$ and fixed $T$
can satisfy
\begin{equation}\label{eq:impossible-observation}
 \norm h_{\HH}^2\le aD_R(h)+b\norm{\mathbf1_{[-T,T]}h}_{\HH}^2
 \quad(h\in\DD\cap\NN^\perp).
\end{equation}
In particular $D_R$ has no positive uniform lower bound by
$\norm h_{\HH}^2$ on this space.
\end{corollary}

The theorem imposes all conditions from the \emph{known source space}
$\NN$. Although $\NN\subset\ker L$ is proved in
Proposition~\ref{prop:source-space}, we do not prove $\NN=\ker L$ and
do not describe the result as removal of the full nullspace. The critical
inequality $\EE(h)\ge\norm h_{\HH}^2/2$ on $\DD\cap\NN^\perp$
remains open here. No operator-norm Weyl sequence or essential-spectrum
conclusion is asserted.

The ingredients are classical theta identities, the compact-test explicit
formula, unconditional zero counting, and divided differences. We specify
their use rather than make a claim of bibliographical priority; a systematic
novelty survey is outside this draft. The constants and starting indices
are proved to exist, but are not numerically certified.

\section{Theta normalization and the source identity}\label{sec:source}

\subsection{The exact theta source}
\begin{lemma}\label{lem:theta}
The function $\kappa$ is positive, smooth, and even. With
$\Xi(z)=\xi(1/2+iz)$ and
$\xi(s)=\tfrac12s(s-1)\pi^{-s/2}\Gamma(s/2)\zeta(s)$, one has
\begin{equation}\label{eq:theta-transform}
 F_\kappa(z):=\int_\R\kappa(x)e^{-izx}\dd x=\Xi(z),
 \qquad
 \int_\R\kappa(x)e^{\pm x/2}\dd x=\frac12.
\end{equation}
For $x\ge0$ there are fixed positive constants $a_0,a_1$ such that
\begin{equation}\label{eq:theta-profile}
 a_0e^{9x/2-\pi e^{2x}}\le\kappa(x)
       \le a_1e^{9x/2-\pi e^{2x}},\qquad a_0=2\pi^2.
\end{equation}
For each fixed integer $r\ge0$,
\begin{equation}\label{eq:theta-derivative}
 \abs{\kappa^{(r)}(x)}\le C_r e^{2r\abs{x}}\kappa(x).
\end{equation}
All fixed derivatives decrease faster than every exponential at both ends.
\end{lemma}
\begin{proof}
Let $S(x)=e^{x/2}\sum_{n\ge1}e^{-\pi n^2e^{2x}}$. The Jacobi
transformation \cite[Eq.~20.7.32]{dlmf} gives
\[
 S(x)-S(-x)=\frac12(e^{-x/2}-e^{x/2}).
\]
Since $\kappa=(\partial_x^2-1/4)S$, the source is even. The series and
each fixed derivative converge locally uniformly. For $x\ge0$,
$z_n\ge\pi>3$, so every original summand is positive.
The Mellin calculation, initially absolutely convergent for $\Re s>1$,
is
\begin{align}\label{eq:mellin}
 \int_\R\kappa(x)e^{(s-1/2)x}\dd x
 &=\zeta(s)\int_0^\infty
       (4\pi^2y^4-6\pi y^2)e^{-\pi y^2}y^{s-1}\dd y\notag\\
 &=\frac12s(s-1)\pi^{-s/2}\Gamma(s/2)\zeta(s).
\end{align}
The standard completed-zeta identities are recorded in
\cite[Eqs.~25.4.3--25.4.4]{dlmf}. The left side continues to an entire
function by the two-sided decay. Evenness and $\xi(0)=\xi(1)=1/2$ give
\eqref{eq:theta-transform} and $\nu(\R)=1$.

For the estimates put $y=\pi e^{2x}$. The first summand and
$4y^2-6y\ge2y^2$ give the lower bound. The upper bound follows from
\[
 \sum_{n\ge1}n^4e^{-(n^2-1)y}
 \le\sum_{n\ge1}n^4e^{-3(n^2-1)}<\infty.
\]
After $r$ derivatives the polynomial in $z_n$ has degree $r+2$:
$P_0(z)=4z^2-6z$ and
$P_{r+1}=P_r/2+2zP_r'-2zP_r$.
The same estimate with $n^{2r+4}$ proves
\eqref{eq:theta-derivative}; its constant may depend on the fixed order.
\end{proof}

\begin{lemma}[Closability]\label{lem:closure}
The form \eqref{eq:energy} is densely defined and closable on $\CC$.
\end{lemma}
\begin{proof}
The combined positive continuous and atomic edge measure is sigma-finite
after restricting bounded endpoint coordinates and shifts bounded away
from zero. Each restricted endpoint marginal is absolutely continuous with
respect to Lebesgue measure; countably many atomic translations preserve
null sets. If $h_k\to0$ in $\HH$ and their edge differences are Cauchy in
edge $L^2$, choose a subsequence converging Lebesgue almost everywhere,
using the strict positivity of the weight. The edge differences then tend
to zero at almost every pair of endpoints. Their edge-$L^2$ limit is zero,
which proves closability. Core energies are finite by the small-shift
derivative bound and theta decay. Core density in $\HH$ follows from the
positive locally smooth weight and even compact approximation.
\end{proof}

\subsection{Classical explicit formula with the constants fixed}
For physical even functions define
\begin{align}\label{eq:Q}
 D_\Gamma(u)&=\int_0^\infty\rho(s)\norm{u-\tau_su}_2^2\dd s,
 \qquad \mathfrak m(u)=\int_\R u(x)c(x)\dd x,\notag\\
 Q(u)&=D_\Gamma(u)-c_\Gamma\norm u_2^2
       -2\sum_{n\ge2}c_n\Re\ip u{\tau_{s_n}u}_2
       +2\abs{\mathfrak m(u)}^2,\notag\\
 c_\Gamma&=\log\pi+\gamma_E+\frac\pi2+3\log2.
\end{align}
Polarization defines $Q(u,v)$. These expressions are not asserted to be
nonnegative.

\begin{lemma}[Classical input in physical coordinates]\label{lem:explicit}
For $u\in C_c^\infty(\R;\C)$ even, let
$F_u(z)=\int_\R u(x)e^{-izx}\dd x$. Then
\begin{equation}\label{eq:explicit}
 Q(u)=\sum_{\rho}
    F_u(z_\rho)\overline{F_u(\overline{z_\rho})},
 \qquad z_\rho=\frac{\rho-1/2}{i},
\end{equation}
where the sum is over nontrivial zeta zeros with multiplicities. The series
converges absolutely. No zero-location hypothesis beyond the classical
critical strip is used.
With the first variable conjugate-linear, the polarized identity is
\begin{equation}\label{eq:explicit-polarized}
 Q(u,v)=\sum_\rho
       \overline{F_u(\overline{z_\rho})}F_v(z_\rho)
 \qquad(u,v\in C_c^\infty(\R;\C)\text{ even}).
\end{equation}
\end{lemma}
\begin{proof}
We apply the classical multiplicative explicit formula for
$\Phi\in C_c^\infty((0,\infty))$, with Mellin transform
$\widetilde\Phi(s)=\int_0^\infty\Phi(y)y^{s-1}\dd y$ and
$\Phi^\sharp(y)=y^{-1}\Phi(y^{-1})$. Its normalization, prime terms,
and coupled archimedean finite part are recorded in
\cite[Appendix B, Eqs.~(147)--(150)]{connesconsani}. Put
\[
 C_u(s)=\int_\R\overline{u(x)}u(x+s)\dd x,
 \qquad \Phi(y)=y^{-1/2}C_u(\log y).
\]
For an even complex $u$, $C_u$ is real and even. Thus $\Phi$ is a permitted
smooth compact multiplicative test. Direct integration gives
\[
 \widetilde\Phi(1/2+iz)
       =F_u(z)\overline{F_u(\bar z)}.
\]
The two pole terms are $2\abs{\mathfrak m(u)}^2$, and the subtracted prime
contribution is $2\sum_{n\ge2}c_n C_u(s_n)$. The negative archimedean
term, after $y=e^s$, is
\begin{equation}\label{eq:archimedean}
 - (\log(4\pi)+\gamma_E)\norm u_2^2
 -2\int_0^\infty
       \frac{e^{s/2}C_u(s)-\norm u_2^2}{e^s-e^{-s}}\dd s.
\end{equation}
The subtraction remains coupled at zero. Using
$\norm{u-\tau_su}_2^2=2(\norm u_2^2-C_u(s))$, this is
$D_\Gamma(u)-c_\Gamma\norm u_2^2$. Indeed the additional constant is
\begin{equation}\label{eq:gamma-constant}
 2\int_0^\infty\frac{e^{s/2}-1}{e^s-e^{-s}}\dd s
   =\psi(1/2)-\psi(1/4)=\frac\pi2+\log2.
\end{equation}
The digamma integral, reflection, and duplication formulas used in this
calculation are \cite[Eqs.~5.9.16, 5.5.4, 5.5.8]{dlmf}.
This proves the precise normalization in \eqref{eq:explicit}.
Compact smooth transforms decay faster than any power on
$\abs{\Im z}\le1/2$. The unconditional bound
$N_\zeta(T)=O(T\log(2+T))$ then proves absolute convergence.
Polarization gives \eqref{eq:explicit-polarized}; conjugation invariance of
the zero multiset makes the displayed pairing Hermitian. More generally,
the mixed series is absolutely convergent whenever both transforms are
$O((1+\abs z)^{-2})$ uniformly on $\abs{\Im z}\le1/2$.
\end{proof}

The only arithmetic zero-counting input used in this paper is the classical
asymptotic
\begin{equation}\label{eq:zero-count}
 N_\zeta(T)=\frac{T}{2\pi}\log\frac{T}{2\pi e}+O(\log T),
\end{equation}
counting multiplicities. For a precise stronger bound see
\cite[Corollary 1.2]{hsw}; none of its explicit numerical constants is used.
The unconditional strip $0<\Re\rho<1$ and its symmetries are recalled in
\cite[Section~25.10]{dlmf}. The factor $1/2$ in our definition of $\xi$
does not affect its zeros.

\subsection{Source cancellation and cutoff continuity}
\begin{lemma}\label{lem:continuity}
On even functions with finite
\begin{equation}\label{eq:Mnorm}
 M(u)=\norm u_{H^1(\R)}+\norm{e^{\abs{x}}u}_2,
\end{equation}
one has $\abs{Q(u,v)}\le C M(u)M(v)$. For every fixed $j\ge0$ and every
compact smooth even $k$,
\begin{equation}\label{eq:radical}
 Q(\kappa^{(2j)},k)=0.
\end{equation}
\end{lemma}
\begin{proof}
The Gamma estimate follows by using
$\norm{u-\tau_su}_2\le s\norm{u'}_2$ for $s\le1$ and
$\norm{u-\tau_su}_2\le2\norm u_2$ for $s\ge1$.
The pole is bounded in $M$ by Cauchy--Schwarz. For the primes,
\begin{equation}\label{eq:correlation}
 \abs{\ip u{\tau_sv}_2}
 \le e^{-\abs{s}}\norm{e^{\abs{x}}u}_2\norm{e^{\abs{x}}v}_2,
 \qquad \sum_{n\ge2}\Lambda(n)n^{-3/2}<\infty.
\end{equation}
This proves the asserted continuity without a sign assumption.

To pass the explicit formula to the source, also use
\[
 S_2(u)=\sum_{r=0}^2\int_\R e^{\abs{x}}\abs{u^{(r)}(x)}\dd x.
\]
Integration by parts gives
$\abs{F_u(z)}\le C S_2(u)(1+\abs z)^{-2}$ on
$\abs{\Im z}\le1/2$. Consequently the polarized zero sum is jointly
continuous under $S_2$ convergence, since
$\sum_\rho(1+\abs{z_\rho})^{-4}<\infty$ by \eqref{eq:zero-count}.
Choose even smooth cutoffs $\chi_R(x)=\chi(x/R)$, with $\chi=1$ on
$[-1,1]$ and supported in $[-2,2]$. For $u=\kappa^{(2j)}$, theta decay
and the product rule imply
\[
 M(\chi_Ru-u)\longrightarrow0,\qquad
 S_2(\chi_Ru-u)\longrightarrow0.
\]
Thus the polarized identity \eqref{eq:explicit-polarized} passes to $(u,k)$.
Since $F_u(z)=(-1)^jz^{2j}\Xi(z)$ vanishes at every zero argument,
the mixed pairing is zero. This proves \eqref{eq:radical}.
\end{proof}

\begin{proposition}[Exact source transform]\label{prop:source-transform}
On $\CC$ one has
\begin{equation}\label{eq:compact-source}
 Q(\kappa h)=\EE(h)-\frac12\Var_\nu(h).
\end{equation}
On the minimum domain $\DD$,
\begin{align}\label{eq:source-identity}
 \EE(h)={}&\frac12\norm h_{\HH}^2+D_\Gamma(\kappa h)
            -c_\Gamma\norm{\kappa h}_2^2\notag\\
 &-2\sum_{n\ge2}c_n\Re\ip{\kappa h}{\tau_{s_n}(\kappa h)}_2.
\end{align}
\end{proposition}
\begin{proof}
For real $a,b$ and complex $u,v$ the edge identity is
\begin{equation}\label{eq:edge-identity}
 \abs{au-bv}^2-(a-b)(a\abs u^2-b\abs v^2)
       =ab\abs{u-v}^2.
\end{equation}
Subtract $Q(\kappa,\kappa\abs h^2)=0$ from $Q(\kappa h)$, using
Lemma~\ref{lem:continuity}. For the continuous part apply
\eqref{eq:edge-identity}; for the primes symmetrize both orientations.
The resulting coefficient of the positive prime-edge integral is exactly
$c_n$, not $2c_n$. The local $c_\Gamma$ terms cancel. Evenness and
\eqref{eq:theta-transform} give
\[
 \mathfrak m(\kappa h)=\frac12\int h\dd\nu,
 \qquad \mathfrak m(\kappa\abs h^2)=\frac12\norm h_{\HH}^2,
 \qquad \mathfrak m(\kappa)=\frac12.
\]
Hence the pole difference is $-\Var_\nu(h)/2$, proving
\eqref{eq:compact-source}. Shift truncation, the coupled differences, and
the continuity bounds above justify the expansion. Substitution of
\eqref{eq:Q} gives \eqref{eq:source-identity} on the core.

For the extension let $Uh=\sqrt{2\kappa c}\,h$ and
$b=\sqrt{\kappa/(2c)}$, so that $\kappa h=bUh$. The map $U$ is unitary,
whereas multiplication by $b$ is merely bounded. Theta decay gives
$b(x)\le C e^{-\abs{x}}$, and therefore
\begin{equation}\label{eq:Kp}
 K_p=\sum_{n\ge2}c_n M_b(\tau_{s_n}+\tau_{-s_n})M_b
\end{equation}
converges absolutely in operator norm:
$\norm{M_b\tau_sM_b}\le C^2e^{-\abs s}$ and
\eqref{eq:correlation} supplies summability. On the core the prime term is
$\ip{Uh}{K_pUh}_2$. The form $D_\Gamma$ is closed by Plancherel, with
nonnegative multiplier
$2\int_0^\infty\rho(s)(1-\cos(ts))\dd s$.

These bounds also give a second proof of closability directly from the
\emph{core} identity \eqref{eq:source-identity}. If $h_k\in\CC$ tends to
zero in $\HH$ and is $\EE$-Cauchy, then for $d=h_k-h_\ell$ the bounded
lower-order terms give
\[
 D_\Gamma(\kappa d)\le\EE(d)+C\norm d_{\HH}^2.
\]
Thus $\kappa h_k=bUh_k$ is Cauchy in the $D_\Gamma$ graph norm and has
$L^2$ limit zero. Closedness yields $D_\Gamma(\kappa h_k)\to0$; the
core identity and the same bounded terms yield $\EE(h_k)\to0$.
This argument uses no prior extension of the identity to $\DD$.

Applied to differences of form-Cauchy core approximants to an arbitrary
$h\in\DD$, the same estimate gives convergence of their physical images
in the $D_\Gamma$ graph norm.
All terms thus pass to the $\HH$ limit. This proves the identity on the
specified minimum closure.
\end{proof}

\begin{lemma}[A minimum-domain admission criterion]\label{lem:admission}
Suppose $h$ is smooth and even, $h\in\HH$, and $u=\kappa h$ and $u'$
decrease faster than every exponential. Then $h\in\DD$ and
\eqref{eq:source-identity} applies to $h$.
\end{lemma}
\begin{proof}
For the cutoffs used above, $\chi_Rh\to h$ in $\HH$ and
$\chi_Ru\to u$ in $M$. For $d=(\chi_R-\chi_S)h\in\CC$, the
\emph{compact} identity gives
\begin{equation}\label{eq:cutoff-difference}
 0\le\EE(d)=Q(\kappa d)+\frac12\Var_\nu(d)
       \le C M(\kappa d)^2+\frac12\norm d_{\HH}^2\longrightarrow0.
\end{equation}
The compact approximants are Cauchy in \eqref{eq:form-norm}; their
$\HH$ limit therefore belongs to the minimum domain. This criterion does
not replace minimum admission by finite maximal energy.
\end{proof}

\begin{proposition}[The known source space]\label{prop:source-space}
For every fixed $j\ge0$, $f_j\in\Dom(L)$, $Lf_j=0$, and
$\ip1{f_j}_{\HH}=4^{-j}$. Consequently $\NN\subset\ker L$.
\end{proposition}
\begin{proof}
Equation~\eqref{eq:theta-derivative} gives
$\abs{f_j(x)}\le C_j e^{4j\abs{x}}$, so $f_j\in\HH$.
Lemma~\ref{lem:admission} applies because $\kappa f_j=\kappa^{(2j)}$.
Polarizing \eqref{eq:compact-source} and using cutoffs and
\eqref{eq:radical}, the form of $L$ satisfies
$\ell(f_j,g)=0$ for every core $g$. Continuity in the form norm extends
this to all $g\in\DD$. The closed-form representation theorem gives
$f_j\in\Dom(L)$ and $Lf_j=0$. Integration by parts yields
\[
 \ip1{f_j}_{\HH}=2\int\kappa^{(2j)}(x)c(x)\dd x
       =4^{-j}2\int\kappa(x)c(x)\dd x=4^{-j}.
\]
The kernel of the closed operator $L$ is $\HH$-closed, so it contains the
closed span \eqref{eq:N}. This proves an inclusion, not equality.
\end{proof}

\section{Exact constraints and clustered zero modes}\label{sec:modes}

Let $Z$ be the zero multiset of $\Xi$. For a zero $z$ of multiplicity $d$
and $0\le r<d$, define
\begin{equation}\label{eq:root-mode}
 v_{z,r}(x)=\frac{\partial_z^r\cos(zx)}{c(x)}.
\end{equation}
\begin{lemma}\label{lem:root-modes}
Every $v_{z,r}$ belongs to $\DD\cap\NN^\perp$. Its real and imaginary
parts have the same membership.
\end{lemma}
\begin{proof}
The functions are even, smooth, and have at most polynomial times
exponential growth. Theta decay gives their $\HH$ membership and the
hypotheses of Lemma~\ref{lem:admission}. For every $j\ge0$,
\begin{align}\label{eq:all-moments}
 \ip{f_j}{v_{z,r}}_{\HH}
 &=2\partial_z^r\int_\R\kappa^{(2j)}(x)\cos(zx)\dd x\notag\\
 &=2\partial_z^r\bigl[(-1)^jz^{2j}\Xi(z)\bigr]=0.
\end{align}
All differentiations and integrations by parts are justified by theta
decay, for each fixed order. The last equality uses $r<d$. Thus every
generator moment, including the mean, vanishes exactly. Closed span gives
$\NN^\perp$; the generators are real, so real and imaginary parts qualify.
\end{proof}

\begin{lemma}[Clusters from zero counting]\label{lem:clusters}
For every sufficiently large integer $m$, there are $m+1$ zero arguments
$z_0,\ldots,z_m$, counted with multiplicity, and a real $\lambda_m$ such
that
\begin{equation}\label{eq:cluster}
 \begin{gathered}
 T_m=e^{m^3},\qquad T_m\le\lambda_m\le2T_m,\\
 \lambda_m\le\Re z_j\le\lambda_m+1/m,
 \qquad \abs{\Im z_j}<1/2.
 \end{gathered}
\end{equation}
\end{lemma}
\begin{proof}
Equation~\eqref{eq:zero-count} gives
$N_\zeta(2T_m)-N_\zeta(T_m)\sim T_m m^3/(2\pi)$.
Partition $[T_m,2T_m]$ into at most $mT_m+1$ intervals of length at most
$1/m$. The average number of zeros is asymptotic to $m^2/(2\pi)$,
which exceeds $m$ for large $m$. Choose an interval containing at least
$m+1$ zeros and let $\lambda_m$ be its left endpoint. The classical strip
gives the imaginary-part bound. An endpoint convention does not affect the
argument. Repeated nodes are taken no more often than their multiplicities.
\end{proof}

No zero at height $T_m$ is computed. The choice $e^{m^3}$ keeps
$\log\lambda_m$ polynomial in $m$ while making $\lambda_m$ dominate the
oscillatory error below.

Write $[z_0,\ldots,z_m]F$ for the divided difference, with its confluent
interpretation at repeated nodes, and define
\begin{equation}\label{eq:P}
 P_m(x)=m!i^{-m}[z_0,\ldots,z_m]\cos(zx),\qquad
 p_m=P_m/c.
\end{equation}
This is a finite linear combination of the modes in
\eqref{eq:root-mode}; therefore $p_m\in\DD\cap\NN^\perp$ exactly.
For distinct nodes, for example,
\[
 P_m(x)=m!i^{-m}\sum_{j=0}^m
          \frac{\cos(z_jx)}{\prod_{k\ne j}(z_j-z_k)}.
\]
Estimating these individual coefficients would be inappropriate when nodes
are close. We instead use the simplex representation.

Let $\theta$ have uniform probability distribution on the simplex
$\theta_j\ge0$, $\sum_{j=0}^m\theta_j=1$, and put
\begin{equation}\label{eq:A}
 A_m(x)=\Etheta\exp\left(ix\sum_{j=0}^m\theta_j(z_j-\lambda_m)\right).
\end{equation}
The classical divided-difference integral formula, discussed in
\cite{deboor}, is
\begin{equation}\label{eq:HG}
 [z_0,\ldots,z_m]F
       =\frac1{m!}\Etheta F^{(m)}\left(\sum_j\theta_jz_j\right).
\end{equation}
It holds for entire $F$ and complex nodes: repeated integration along
complex straight segments proves the formula for distinct nodes, and
continuous extension gives the confluent case. Equivalently, its power
series proof follows from the simplex moments and complete homogeneous
polynomials. Applying it to the two exponentials in the cosine gives
\begin{equation}\label{eq:P-simplex}
 P_m(x)=\frac{x^m}{2}\left(
       e^{i\lambda_mx}A_m(x)+(-1)^m e^{-i\lambda_mx}A_m(-x)\right).
\end{equation}

\begin{lemma}\label{lem:amplitudes}
For $m\ge2$ and real $x$,
\begin{align}\label{eq:A-bounds}
 \abs{A_m(\pm x)},\ \abs{A_m'(\pm x)}&\le e^{\abs{x}/2},\notag\\
 \abs{P_m(x)}&\le\abs{x}^m e^{\abs{x}/2},\notag\\
 \abs{P_m'(x)}&\le
       \left(m\abs{x}^{m-1}+(\lambda_m+1)\abs{x}^m\right)e^{\abs{x}/2}.
\end{align}
If $0\le x\le b_m:=m/4$, then
\begin{equation}\label{eq:A-lower}
 \Re A_m(\pm x)\ge\cos(1/4)e^{-x/2}.
\end{equation}
\end{lemma}
\begin{proof}
Each convex combination of $z_j-\lambda_m$ has real part in $[0,1/m]$
and imaginary part in $(-1/2,1/2)$; its modulus is less than one for
$m\ge2$. This proves the first line and its derivative bound directly
from \eqref{eq:A}, and the next two lines follow from
\eqref{eq:P-simplex}. On the indicated interval the phase magnitude is
at most $x/m\le1/4$, whereas the positive exponential amplitude is at
least $e^{-x/2}$. This proves \eqref{eq:A-lower}.
\end{proof}

\section{Normalization and separation of theta moments}\label{sec:moments}

Put $W=\kappa/c$ on $[0,\infty)$ and define
\begin{equation}\label{eq:d}
 d_m=\int_0^{b_m}x^{2m}W(x)e^{-x}\dd x.
\end{equation}
This is a concrete positive source moment for the constructed functions,
not an unknown optimal constant for the critical form inequality.

\begin{lemma}[Control of the oscillatory cross term]\label{lem:norm}
For all sufficiently large $m$,
\begin{equation}\label{eq:norm-lower}
 \norm{p_m}_{\HH}^2\ge d_m/2.
\end{equation}
\end{lemma}
\begin{proof}
Expand \eqref{eq:P-simplex}. The diagonal terms and
\eqref{eq:A-lower} give
\begin{align}\label{eq:cross}
 \int_0^{b_m}W\abs{P_m}^2
 &\ge\frac{\cos^2(1/4)}2d_m-\frac12\abs{I_m},\notag\\
 I_m&=\int_0^{b_m}x^{2m}W(x)A_m(x)\overline{A_m(-x)}
                                  e^{2i\lambda_mx}\dd x.
\end{align}
Let $B_m(x)=x^{2m}W(x)A_m(x)\overline{A_m(-x)}$. Then $B_m(0)=0$,
and integration by parts yields
\begin{equation}\label{eq:cross-bound}
 \abs{I_m}\le
 \frac{\abs{B_m(b_m)}+\int_0^{b_m}\abs{B_m'}\dd x}{2\lambda_m}
 \le\frac{C(m+1)(1+m)^{2m}}{\lambda_m}.
\end{equation}
Indeed \eqref{eq:A-bounds} bounds the integrand in the numerator by
\[
 2m x^{2m-1}W(x)e^x+x^{2m}(|W'(x)|+2W(x))e^x.
\]
The fixed source weights have finite integrals and suprema by
Lemma~\ref{lem:theta}. For $m\ge8$,
$d_m\ge d_0:=\int_1^2W(x)e^{-x}\dd x>0$.
Since $\lambda_m\ge e^{m^3}$, the right side of
\eqref{eq:cross-bound} is $o(d_m)$. Finally, evenness and the exact
weight give
\[
 \norm{p_m}_{\HH}^2=4\int_0^\infty W(x)\abs{P_m(x)}^2\dd x.
\]
Together with \eqref{eq:cross} this proves \eqref{eq:norm-lower}.
\end{proof}

For a fixed $a\ge0$ write
\begin{equation}\label{eq:U}
 U_m(a)=\int_0^\infty x^{2m}\frac{\kappa(x)^2}{c(x)^2}e^{ax}\dd x,
 \qquad M_m=\frac{\log m}{4},\qquad C_0=\pi e^4.
\end{equation}
\begin{lemma}[Separation of the two theta weights]\label{lem:moment-ratio}
For every fixed $a\ge0$ there is $C_a<\infty$ such that, for all
sufficiently large $m$,
\begin{equation}\label{eq:moment-ratio}
 \frac{U_m(a)}{d_m}\le C_a e^{-\sqrt m},\qquad
 \frac1{d_m}\le C e^{-\sqrt m}.
\end{equation}
The starting index may depend on $a$.
\end{lemma}
\begin{proof}
For large $m$, $[M_m+1,M_m+2]\subset[0,b_m]$. Since
$e^{2M_m}=\sqrt m$ and $c(x)\le e^{x/2}$, the lower theta estimate gives
\begin{equation}\label{eq:d-lower}
 d_m\ge2\pi^2(M_m+1)^{2m}e^{-C_0\sqrt m}.
\end{equation}
In particular $d_m^{-1}=O(e^{-\sqrt m})$.

Split $U_m(a)$ at $M_m$ and $b_m$. On $[0,M_m]$ it is bounded by
$C_a M_m^{2m}$, since its fixed source weight is integrable. Consequently
\begin{equation}\label{eq:ratio-low}
 \frac{U_{m,\mathrm{low}}(a)}{d_m}
 \le C_a\exp\left(-\frac{2m}{M_m+1}+C_0\sqrt m\right).
\end{equation}
Here $\log(M/(M+1))\le-1/(M+1)$, and $m/\log m$ dominates
$\sqrt m$, giving the required rate for this part.

On $[M_m,b_m]$ the ratio of the numerator weight to that in $d_m$ is
\begin{equation}\label{eq:ratio-mid}
 \frac{\kappa(x)}{c(x)}e^{(a+1)x}
 \le C_a e^{-(\pi/2)e^{2x}}
 \le C_a e^{-(\pi/2)\sqrt m}.
\end{equation}
This follows from \eqref{eq:theta-profile} by absorbing a fixed
exponential factor into half of the theta exponent. Integrating bounds
the middle contribution by the last quantity times $d_m$.

On $[b_m,\infty)$ the source weight is at most $C_a e^{-\pi e^{2x}}$.
For $x\ge m/4$ and all sufficiently large $m$,
$2m\log x\le(\pi/2)e^{2x}$ uniformly. It suffices to check the ratio
at $m/4$ and use the subsequent decrease of $\log x/e^{2x}$.
The tail is therefore bounded by $C_a e^{-c e^{m/2}}$ for a fixed $c>0$.
Divide by $d_m\ge d_0>0$ and combine the three parts to obtain
\eqref{eq:moment-ratio}.
\end{proof}

Splitting at $x=1$ also gives
\begin{equation}\label{eq:previous-moment}
 U_{m-1}(1)\le U_m(1)+C.
\end{equation}
Thus any fixed polynomial in $m$ times
$(U_m(a)+U_{m-1}(1)+1)/d_m$ tends to zero. The estimates are asymptotic
analytic inequalities, not finite numerical enclosures.

\section{Energy control and escape from the center}\label{sec:energy}

Let $u_m=\kappa p_m=\kappa P_m/c$. Lemma~\ref{lem:amplitudes} and
$\abs{(\kappa/c)'}\le C e^{2x}\kappa/c$ for $x\ge0$ imply
\begin{align}\label{eq:physical-bounds}
 \norm{u_m}_2^2&\le2U_m(1),\notag\\
 \norm{e^{\abs{x}}u_m}_2^2&\le2U_m(3),\notag\\
 \norm{u_m'}_2^2&\le C\bigl((\lambda_m+1)^2U_m(1)
                       +U_m(5)+m^2U_{m-1}(1)\bigr).
\end{align}
These constants are independent of the chosen cluster and of $m$.

\begin{lemma}[Logarithmic Gamma estimate]\label{lem:gamma-log}
For $u\in H^1(\R)$ and $\Omega\ge2$,
\begin{equation}\label{eq:gamma-log}
 D_\Gamma(u)\le C\left(\Omega^{-2}\norm{u'}_2^2
                   +(1+\log\Omega)\norm u_2^2\right).
\end{equation}
\end{lemma}
\begin{proof}
Split the shift integral at $1/\Omega$ and $1$. On $(0,1/\Omega)$ use
$\norm{u-\tau_su}_2\le s\norm{u'}_2$ and $\rho(s)\le C/s$.
On $(1/\Omega,1)$ use $2\norm u_2$, giving $C\log\Omega$.
On $(1,\infty)$ the kernel is integrable. This proves the bound directly
from the shift integral, without a separate multiplier asymptotic.
\end{proof}

With $\Omega=\lambda_m+2$ one has $\log\Omega\le m^3+C$. Equations
\eqref{eq:moment-ratio}, \eqref{eq:previous-moment}, and
\eqref{eq:physical-bounds} give
\begin{equation}\label{eq:physical-small}
 D_\Gamma(u_m)+\norm{e^{\abs{x}}u_m}_2^2
       \le C m^3e^{-\sqrt m}d_m.
\end{equation}
The prime term retains its arithmetic weights:
\begin{equation}\label{eq:prime-small}
 2\sum_{n\ge2}c_n\abs{\ip{u_m}{\tau_{s_n}u_m}_2}
 \le 2\sum_{n\ge2}\frac{\Lambda(n)}{n^{3/2}}
               \norm{e^{\abs{x}}u_m}_2^2.
\end{equation}
Using the exact identity \eqref{eq:source-identity}, we obtain
\begin{equation}\label{eq:unnormalized-error}
 \abs{\EE(p_m)-\tfrac12\norm{p_m}_{\HH}^2}
        \le C m^3e^{-\sqrt m}d_m.
\end{equation}

Choose $r_m$ to be the real or imaginary part of $p_m$ having the larger
$\HH$ norm, and put
\begin{equation}\label{eq:normalized}
 h_m=\frac{r_m}{\norm{r_m}_{\HH}},\qquad
 \norm{r_m}_{\HH}^2\ge\frac12\norm{p_m}_{\HH}^2\ge\frac14d_m.
\end{equation}
These are real smooth even members of $\DD\cap\NN^\perp$.
The estimates \eqref{eq:physical-bounds}--\eqref{eq:prime-small} hold for
each real or imaginary part with the same upper bounds. Applying the
source identity to the chosen part separately proves
\eqref{eq:main-energy} after division by \eqref{eq:normalized}. This step
does not infer a signed bound for one part from possible cancellation
between the two signed errors.

For fixed $T>0$, \eqref{eq:A-bounds}, \eqref{eq:d-lower}, and
\eqref{eq:normalized} imply
\begin{equation}\label{eq:local-explicit}
 \norm{\mathbf1_{[-T,T]}h_m}_{\HH}^2
 \le C_T\left(\frac{T}{M_m+1}\right)^{2m}e^{C_0\sqrt m}\longrightarrow0.
\end{equation}
The norm-one bound and compact truncation of arbitrary $\HH$ tests then
give weak convergence to zero. The same calculation with $T=M_m$ uses
$\int_0^\infty W(x)e^x\dd x<\infty$ and yields the uniform moving-cutoff
bound
\begin{equation}\label{eq:moving-cutoff}
 \norm{\mathbf1_{[-M_m,M_m]}h_m}_{\HH}^2
 \le C\exp\left(-\frac{2m}{M_m+1}+C_0\sqrt m\right).
\end{equation}
This is $O(e^{-\sqrt m})$ after increasing the starting index.
It completes the proof of \eqref{eq:main-energy}--\eqref{eq:main-local}.

\section{A two-shift envelope and its minimum-domain transform}\label{sec:envelope}

\subsection{Safe shifts and existence of a positive envelope}
For $\abs t<a_*:=(\log2)/2$ define
\begin{equation}\label{eq:shifts}
 \Phi_t(x)=\frac{\kappa(x+t)+\kappa(x-t)}{2\kappa(x)},
 \qquad q_t=\Phi_t-\cosh(t/2).
\end{equation}
\begin{lemma}\label{lem:shifts}
For every safe real shift, $q_t\in\NN\cap\Dom(A)$,
$\ip1{q_t}_{\HH}=0$, and $Aq_t=q_t/2$.
\end{lemma}
\begin{proof}
The theta series is holomorphic for $\abs{\Im x}<\pi/4$. On compact
subsets of the connected region
\begin{equation}\label{eq:shift-region}
 \mathcal U=\left\{t\in\C:\abs{\Im t}<\pi/4,
         \ e^{-2\abs{\Re t}}\cos(2\Im t)>1/2\right\},
\end{equation}
the shifted source, squared and divided by $\kappa$, is bounded in the
tails by a fixed exponential times
\[
 \exp\left(-\pi\bigl[2e^{-2\abs{\Re t}}\cos(2\Im t)-1\bigr]
                 e^{2\abs x}\right).
\]
The bracket has a positive lower bound there. Local domination, also for
local $t$ derivatives, proves that $q_t$ is $\HH$-holomorphic.
Near zero its Taylor series is
\[
 q_t=\sum_{j\ge1}\frac{t^{2j}}{(2j)!}(f_j-4^{-j}).
\]
The quotient map to $\HH/\NN$ is therefore zero near zero, and hence on
$\mathcal U$ by analytic continuation. This proves $q_t\in\NN$ on the
whole safe real interval without assuming a global Taylor radius.
Changing variables gives $\int\Phi_t\dd\nu=\cosh(t/2)$.
Proposition~\ref{prop:source-space} now gives the action and domain claims.
\end{proof}

\begin{definition}\label{def:envelope}
An admissible envelope has parameters satisfying
\begin{equation}\label{eq:envelope-parameters}
 0<t_2<t_1<a_*,\qquad
 2e^{-2t_2}>1+e^{-2t_1},\qquad w_1,w_2>0,
\end{equation}
such that $q_{t_1},q_{t_2}$ have no common real zero. Set
\begin{equation}\label{eq:envelope}
 R=\sqrt{w_1q_{t_1}^2+w_2q_{t_2}^2},\qquad
 \Delta=R-\sqrt{w_1}q_{t_1},\qquad K=\kappa R.
\end{equation}
\end{definition}

\begin{lemma}\label{lem:envelope-exists}
For each fixed $t_1\in(0,a_*)$ an admissible $t_2$ exists. For any
admissible parameters, $R>0$, $\inf R>0$, and $R$ tends to infinity
at both ends. The function $\Delta\ge0$ and all its fixed derivatives
decrease faster than every exponential. Moreover
\begin{equation}\label{eq:gamma-range}
 1<\gamma:=\sup_x\frac{\Delta(x)}{R(x)}\le2.
\end{equation}
\end{lemma}
\begin{proof}
For fixed safe $t>0$, the leading theta term gives
\begin{equation}\label{eq:q-asymptotic}
 q_t(x)\sim\frac12e^{-9t/2}
       \exp\left(\pi(1-e^{-2t})e^{2\abs x}\right).
\end{equation}
Thus $q_{t_1}$ is positive in both tails and, being nonzero and real
analytic, has finitely many real zeros. Write $a_i=e^{-2t_i}$ and
$t_*=-\tfrac12\log((1+a_1)/2)$, so $0<t_*<t_1$.
At each zero $x$ of $q_{t_1}$, the real analytic function $t\mapsto q_t(x)$
is not identically zero: as $t\to+\infty$ its shifted numerator tends to
zero while $\cosh(t/2)$ diverges. Its zeros on $[0,t_*]$ are finite.
Choose $t_2\in(0,t_*)$ outside the finite union of these exceptions.
This proves existence, without certifying a numerical pair.

In the tails $q_{t_1}>0$ and $q_{t_2}/q_{t_1}\to0$, giving
\begin{equation}\label{eq:Delta-asymptotic}
 R\sim\sqrt{w_1}q_{t_1},\qquad
 \Delta\sim\frac{w_2}{2\sqrt{w_1}}\frac{q_{t_2}^2}{q_{t_1}}.
\end{equation}
The latter follows from
$\Delta=w_2q_{t_2}^2/(R+\sqrt{w_1}q_{t_1})$ \emph{in the tails};
this quotient need not have a nonzero denominator everywhere.
Its exponential coefficient is $-\eta$, where
$\eta=\pi(2a_2-1-a_1)>0$.
Fixed differentiation of the theta series and the quotient in the tails
therefore gives
\begin{equation}\label{eq:Delta-decay}
 \abs{\Delta^{(r)}(x)}\le C_r e^{C_r\abs x}e^{-\eta e^{2\abs x}}
\end{equation}
for large $\abs x$, weakening $\eta$ if necessary. To justify the
quotient derivatives, use
$R=\sqrt{w_1}q_{t_1}(1+(w_2/w_1)(q_{t_2}/q_{t_1})^2)^{1/2}$ there;
fixed logarithmic derivative bounds grow only polynomially in $e^{2\abs x}$.
On the remaining compact interval $R$ has no zero and is smooth.
This proves all the asserted regularity and the positive infimum of $R$.

Finally $\Delta/R=1-\sqrt{w_1}q_{t_1}/R\in[0,2]$. The nonzero
mean-zero function $q_{t_1}$ must be negative somewhere, because it is
positive in the tails. At such a point $\Delta/R>1$, proving
\eqref{eq:gamma-range}.
\end{proof}

\subsection{The bounded potential and the ground-state identity}
\begin{lemma}[A prime-weighted row bound]\label{lem:prime-row}
If $J\ge0$ is continuous and $J(y)\le C_Be^{-B\abs y}$ for some
$B>1/2$, then
\begin{equation}\label{eq:prime-row}
 \sum_{n\ge2}\frac{\Lambda(n)}{\sqrt n}
       \bigl(J(x+s_n)+J(x-s_n)\bigr)\le C_J c(x).
\end{equation}
\end{lemma}
\begin{proof}
For positive integer $r$, the prime-adic valuations of the central
binomial coefficient give
$\psi(2r)-\psi(r)\le\log\binom{2r}{r}\le2r\log2$.
Dyadic summation and monotonicity imply the elementary Chebyshev estimate
$\psi(X)\le CX$. Thus
\[
 \sum_{k\le\log n<k+1}\frac{\Lambda(n)}{\sqrt n}\le Ce^{k/2}.
\]
Multiply by the supremum of $J(x\pm s)$ on each unit interval. The
inequality $k\le\abs x+\abs{x\pm k}$ bounds the result by $Ce^{\abs x/2}$
times a uniformly summable lattice series with decay exponent $B-1/2$.
Since $e^{\abs x/2}\le2c(x)$, this proves the claim. The weights have not
been discarded.
\end{proof}

\begin{proposition}[Minimum-domain ground-state transform]\label{prop:ground}
For an admissible envelope, $\Delta,R\in\Dom(A)$ and $A\Delta$ is
bounded. The function
\begin{equation}\label{eq:V}
 V=\frac12-\frac{AR}{R}=\frac{\Delta/2-A\Delta}{R}
\end{equation}
is bounded and, for some $b_R>0$,
\begin{equation}\label{eq:V-decay}
 \abs{V(x)}\le C_R e^{-b_Re^{2\abs x}}
 \quad\text{for large }\abs x.
\end{equation}
For every $h\in\DD$ the following transformed energy is finite and
satisfies
\begin{align}\label{eq:ground}
 D_R(h)&=\iint K(x)K(z)
       \abs{\frac{h(x)}{R(x)}-\frac{h(z)}{R(z)}}^2
                                  \dd\mu(x-z)\dd x,\notag\\
 \EE(h)-\frac12\norm h_{\HH}^2
       &=D_R(h)-\int_\R V\abs h^2\dd\nu.
\end{align}
\end{proposition}
\begin{proof}
Equation~\eqref{eq:Delta-decay} and Lemma~\ref{lem:admission} give
$\Delta\in\DD$. Define the coupled expression
\begin{equation}\label{eq:A-Delta}
 B_\Delta(x)=\frac1{c(x)}\int\kappa(z)
                  (\Delta(x)-\Delta(z))\dd\mu(x-z).
\end{equation}
For shifts of size less than one use the bounded derivative of
$\Delta$ and $\int_0^1s\rho(s)\dd s<\infty$; do not separate the
two singular convolutions. The large continuous shifts have finite mass.
For the primes use Lemma~\ref{lem:prime-row} with $J=\kappa$, and the
boundedness of $\Delta$. These estimates prove
$\norm{B_\Delta}_\infty<\infty$.
Symmetrization of the half measure against a compact test yields
$\EE(\Delta,\varphi)=\ip{B_\Delta}\varphi_{\HH}$.
Density in $\DD$ and the representation theorem give
$\Delta\in\Dom(A)$ and $A\Delta=B_\Delta$.

Lemma~\ref{lem:shifts} and $R=\sqrt{w_1}q_{t_1}+\Delta$ now give
$R\in\Dom(A)$ and
$AR=\sqrt{w_1}q_{t_1}/2+A\Delta$. This proves \eqref{eq:V}.
The positive infimum and tail asymptotic of $R$, together with the
bounded numerator in \eqref{eq:V}, prove \eqref{eq:V-decay} and boundedness.
No global bound for $R'$ is required.

For $h\in\CC$, $g=h/R$ and $\abs h^2/R$ are compact smooth functions.
Applying \eqref{eq:edge-identity} with $a=R(x)$ and $b=R(z)$ gives
\[
 D_R(h)=\EE(h)-\EE(R,\abs h^2/R)
        =\EE(h)-\int\frac{AR}{R}\abs h^2\dd\nu.
\]
Shift truncation, edge Cauchy--Schwarz for the mixed term, and monotone
convergence for the positive term justify the integration. This is
\eqref{eq:ground} on the core.

For $h\in\DD$, let core functions $h_k\to h$ in the form norm. Boundedness
of $V$ and the identity applied to differences show that $h_k/R$ is Cauchy
in the positive transformed edge energy. It also converges to $h/R$ in
$L^2(R^2\dd\nu)$. A subsequence converges at both endpoints almost
everywhere for the sigma-finite transformed edge measure, because its
restricted endpoint marginals are Lebesgue absolutely continuous, including
every countable atomic translation. The edge-$L^2$ limit is thus the actual
difference of $h/R$. Passing to the limit proves both finiteness and the
identity on the minimum domain. No maximal-domain equality has been used.
\end{proof}

\subsection{Completion of the escape theorem and its boundary}
Split the potential integral at $M_m=(\log m)/4$. Equations
\eqref{eq:moving-cutoff} and \eqref{eq:V-decay} imply, for some $a_R>0$,
\begin{equation}\label{eq:potential-small}
 \int\abs V\abs{h_m}^2\dd\nu
 \le \norm V_\infty\norm{\mathbf1_{[-M_m,M_m]}h_m}_{\HH}^2
       +\sup_{\abs x>M_m}\abs{V(x)}
 \le C_Re^{-a_R\sqrt m}.
\end{equation}
The ground-state identity and \eqref{eq:main-energy} prove
\eqref{eq:main-ground}, after decreasing $a_R$ if necessary. This completes
Theorem~\ref{thm:escape}. Substituting this sequence proves
Corollary~\ref{cor:obstructions}.

For real $h=Rg\in\DD\cap\NN^\perp$, put
\begin{equation}\label{eq:critical}
 F(g)=\int g^2R^2V\dd\nu,\qquad D(g)=D_R(Rg).
\end{equation}
The desired critical comparison is $F(g)\le D(g)$. Since
$\NN\subset\ker L$, the space $\NN$ reduces $L$, and decomposing a
form vector into its $\NN$ and $\NN^\perp$ parts shows that this
comparison is equivalent to $L\ge0$. This is an equivalence, not an
additional estimate. On the constructed sequence both $F(h_m/R)$ and
$D(h_m/R)$ tend to zero; their relative ordering is not determined.

\section{Scope, verification, and data availability}\label{sec:limits}

The functions in Theorem~\ref{thm:escape} are noncompact. They satisfy
all conditions from $\NN$ exactly and belong to the prescribed minimum
form domain. They are not compact tail bumps with only finitely many
moments removed. Their small mass on every fixed center does not assert
exact vanishing on such a center. The bound is stated for sufficiently
large $m$; neither a finite starting index nor numerical values of its
constants have been certified. The cluster height $e^{m^3}$ is an
analytic existence device, not a performed or proposed zero computation.

The critical sign in \eqref{eq:critical} and the Riemann hypothesis remain
open in this work. We do not prove $\NN=\ker L$, global completeness of
the zero modes in $\NN^\perp$, membership of the entire sequence in
$\Dom(L)$, or $\norm{Lh_m}\to0$. In particular, convergence of form
Rayleigh quotients is not presented as an operator Weyl-sequence argument
or an essential-spectrum conclusion.

The accompanying research artifacts include an exact integer-coefficient
check of \eqref{eq:edge-identity} and of the Young square used in
Appendix~\ref{app:signed}. These checks verify polynomial identities only;
they do not verify the explicit formula, analytic estimates, domain
arguments, or zero counting. No theta, zero, physical integral, FFT, Gram,
or eigenvalue computation is used to establish these results. Compilation and
rendering checks concern document integrity and layout.

The proof notes, source and dependency audits, reproduction instructions,
and saved exact algebra records are preserved at the following fixed
repository commit:
\begin{quote}\small
\url{https://github.com/serabing-hash/theta-weighted-certificate/commit/2447014b009ed432b353d6a3817294df59ff712e}.
\end{quote}
This artifact snapshot was publicly accessible before the manuscript's
repository version was created. Its file hashes identify the supporting
material; they are not mathematical proof certificates.

\appendix
\section{Local density and compact-support uniqueness}\label{app:local}

\begin{proposition}\label{prop:local-density}
For each finite $S>0$, the restrictions of the modes
\eqref{eq:root-mode} have dense linear span in
$L^2_{\mathrm{even}}([-S,S],\nu;\C)$.
\end{proposition}
\begin{proof}
If $v$ annihilates all restricted modes, then
\[
 B(z)=2\int_{-S}^S\kappa(x)\overline{v(x)}\cos(zx)\dd x
\]
is entire and $\abs{B(z)}\le C_v e^{S\abs{\Im z}}$. By orthogonality
it vanishes at every zero of $\Xi$ to at least its multiplicity.
A nonzero entire function of this exponential type has $O(r)$ zeros in
$\abs z\le r$, counting multiplicity: remove any finite-order zero at
zero and apply Jensen's formula at $2r$. Equation~\eqref{eq:zero-count}
instead supplies a superlinear count. Thus $B\equiv0$, and Fourier
uniqueness for the compact $L^1$ function $\kappa\overline v$ gives
$v=0$. This proves density. Real and imaginary parts give the real version.
\end{proof}

\begin{proposition}\label{prop:compact-unique}
If $h\in\NN^\perp$ has compact support, then $h=0$.
\end{proposition}
\begin{proof}
The even function $f=hc$ is compact $L^1$, and $\kappa*f(t)$ is
holomorphic on $\abs{\Im t}<\pi/4$ by the theta series and compact
support. At zero, all even derivatives vanish by the generator moments;
all odd derivatives vanish by parity. The identity theorem gives
$\kappa*f=0$ on the strip. Fourier transformation on the real line yields
$\Xi\widehat f=0$. The real zeros of the nonzero entire $\Xi$ are
isolated, so continuity of $\widehat f$ and Fourier uniqueness give $f=0$.
\end{proof}

This argument does not apply to arbitrary unbounded tails, since $hc$
need not be tempered for general $h\in\HH$. Proposition~\ref{prop:local-density}
is only a local restriction statement and does not imply a global
expansion of $\NN^\perp$ in zero modes.

\section{An oriented signed-jump estimate}\label{app:signed}

For real $h\in\CC$ and $g=h/R$, define
\begin{equation}\label{eq:J-signed}
 J(g)=\iint\kappa(x)\kappa(z)(\Delta(z)-\Delta(x))
             (R(x)g(x)^2-R(z)g(z)^2)\dd\mu(x-z)\dd x.
\end{equation}
Let $\Omega_+=\{(x,z):\Delta(z)>\Delta(x)\}$. The integrand is symmetric
under swapping endpoints, so $J$ is twice its integral over $\Omega_+$.
For $\alpha>0$ Young's inequality gives
\[
 R(x)g(x)^2-R(z)g(z)^2
 \le(1+\alpha)R(x)(g(x)-g(z))^2
       +\bigl((1+1/\alpha)R(x)-R(z)\bigr)g(z)^2.
\]
On $\Omega_+$, $\Delta(z)-\Delta(x)\le\Delta(z)\le\gamma R(z)$.
Twice the integral of the symmetric transformed squared difference over
$\Omega_+$ is at most its full integral, including possible equality edges.
Consequently
\begin{align}\label{eq:signed-bound}
 J(g)&\le(1+\alpha)\gamma D(g)+\int_\R g(z)^2\widetilde W(z)\dd z,\notag\\
 \widetilde W(z)&=2\kappa(z)\int_{\Delta(z)>\Delta(x)}
       \kappa(x)(\Delta(z)-\Delta(x))\notag\\
 &\hspace{28mm}\cdot\bigl[(1+1/\alpha)R(x)-R(z)\bigr]_+
                                                  \dd\mu(x-z).
\end{align}
The last integral in the first line is with respect to Lebesgue measure.
For the stated compact tests, the $\Delta$ difference removes the small-shift
singularity; large shifts use the prime-weighted row bound for $K$.
No unrestricted nonlinear form pairing on all of $\DD$ is asserted.

This orientation saves a factor of two in the displayed jump coefficient,
but \eqref{eq:gamma-range} gives
\[
 1+\alpha<(1+\alpha)\gamma\le2(1+\alpha).
\]
Moreover $F(g)=\tfrac12\int\Delta Rg^2\dd\nu+J(g)$ on the core.
Thus \eqref{eq:signed-bound} does not absorb $D$ with a coefficient below
one and does not prove the critical comparison.

\begin{thebibliography}{9}

\bibitem{connesconsani}
A.~Connes and C.~Consani,
\emph{Weil positivity and Trace formula: the archimedean place},
arXiv:2006.13771v1 (2020), Appendix~B, Eqs.~(147)--(150).
\url{https://arxiv.org/abs/2006.13771v1}.

\bibitem{deboor}
C.~de Boor, Divided differences,
\emph{Surveys in Approximation Theory} \textbf{1} (2005), 46--69.
\url{https://arxiv.org/abs/math/0502036}.

\bibitem{hsw}
E.~Hasanalizade, Q.~Shen, and P.-J.~Wong,
\emph{Counting zeros of the Riemann zeta function},
arXiv:2107.06506 (2021), Corollary~1.2.
\url{https://arxiv.org/abs/2107.06506}.

\bibitem{dlmf}
NIST Digital Library of Mathematical Functions,
Version~1.2.8, released 15 September 2026.
Eqs.~20.7.32, 25.4.3--25.4.4, 5.9.16, 5.5.4, 5.5.8;
Section~25.10. Accessed 17 October 2026.
\url{https://dlmf.nist.gov/}.

\end{thebibliography}
\end{document}
